% =====================================================================
% Section 7 — Conclusion (short, RESS-style)
% =====================================================================
% Per professor's recommendation, the conclusion is condensed to a
% restatement of the main findings; the experimental roadmap and the
% candidate barrier extensions have been moved to Section 6
% (Validation roadmap and future work).
% =====================================================================

\section{Conclusion}\label{sec:conclusion}

This work proposed a quantitative-coupled BowTie risk analysis for
the valve train of a high-performance VW AP 1.8 cylinder head
operating at 10\,000~rpm. The qualitative bow diagram was enriched
with five deterministic indicators anchored to literature-based pass
criteria, and the indicators were propagated through a Monte Carlo
refinement under ASTM~A877/A877M wire tolerances and the
strength dispersion of \citet[\S~10-30]{shigley2020}
($N=2\times 10^{5}$ trials). Three findings result: (i)~two
concurrent at-risk preventive barriers --- the outer-spring
coil-bind margin ($M_{\mathrm{out}}=0.78$~mm, $P[\mathrm{fail}]=77.9\,\%$)
and the inner-spring Goodman factor of safety
($n_{f,\mathrm{in}}\approx 0.68$, $P[\mathrm{fail}]=100\,\%$); (ii)~a
residual concern on the outer-spring natural-frequency margin
($f_{n,\mathrm{out}}/f_{\mathrm{cam,max}}\approx 3.81$ against the
conventional $4\times$ rule of thumb); and (iii)~the procedure
identifies multiple concurrent critical barriers that a conventional
AIAG--VDA scored FMEA collapses into a narrow RPN band. The
empirical confirmation of these analytical findings on the AP 1.8
head, together with the extension of the barrier set to assembly
tolerance and lubrication regime motivated by the construct-validity
benchmark, is the subject of the validation roadmap of
Section~\ref{sec:roadmap}.
